\chapter{Class and Class Conflict}
% \addcontentsline{toc}{chapter}{Class and Class Conflict}

\section*{1}
At the core of Marxist politics, there is the notion of conflict. But this is not what makes it specific and distinct: all concepts of politics, of whatever kind, are about conflict - how to contain it, or abolish it. What is specific about Marxist politics is what it declares the nature of the conflict to be; and what it proclaims to be its necessary outcome.

In the liberal view of politics, conflict exists in terms of 'problems' which need to be 'solved'. The hidden assumption is that conflict does not, or need not, run very deep; that it can be 'managed' by the exercise of reason and good will, and a readiness to compromise and agree. On this view, politics is not civil war conducted by other means but a constant process of bargaining and accommodation, on the basis of accepted procedures, and between parties who have decided as a preliminary that they could and wanted to live together more or less harmoniously. Not only is this sort of conflict not injurious to society: it has positive advantages. It is not only civilized, but also civilizing. It is not only a means of resolving problems in a peaceful way, but also of producing new ideas, ensuring progress, achieving ever-greater harmony, and so on. Conflict is 'functional', a stabilizing rather than a disruptive force.

The Marxist approach to conflict is very different. It is not a matter of 'problems' to be 'solved' but of a state of domination and subjection to be ended by a total transformation of the conditions which give rise to it. No doubt conflict may be attenuated, but only because the ruling class is able by one means or another - coercion, concessions, or persuasion - to prevent the subordinate classes from seeking emancipation. Ultimately, stability is not a matter of reason but of force. The antagonists are irreconcilable, and the notion of genuine harmony is a deception or a delusion, at least in relation to class societies.

For the protagonists are not individuals as such, but individuals as members of social aggregates - classes. In the Grund- risse, Marx writes that 'society does not consist of individuals, but expresses the sum of interrelations, the relations within which these individuals stand. As if someone were to say: Seen from the perspective of society, there are no slaves and no citizens: both are human beings. Rather they are that outside society. To be a slave, to be a citizen, are social characteristics, relations between human beings A and B. Human being A, as such, is not a slave. He is a slave in and through society.'\endnote{K. Marx, Grundrisse, p. 265.} A member of one class may well feel no antagonism towards members of other classes; and there may be mobility between classes. But classes nevertheless remain irreconcilably divided, whether conflict occurs or not, and independently of the forms it may or may not assume. It is important not to attribute to the notion of conflict the meaning of 'eruption or of an interruption of an otherwise smooth, harmonious process. This is often the meaning implicit in the liberal usage of 'conflict ' - the expression of a 'problem ' or 'problems ' that need to be 'solved '. In a Marxist perspective, this is a mystification: conflict is inherent in the class system, incapable of solution within that system. Eruptions, outbursts, revolts, revolutions, are only the most visible manifestations of a permanent alienation and conflict, signs that the contradictions in the social system are growing and that the struggle between contending classes is assuming sharper or irrepressible forms. These contending classes are locked in a situation of domination and subjection from which there is no escape except through the total transformation of the mode of production.

Domination is a central concept in Marxist sociology and politics. But domination, in Marxist thought, is not an inherent part of 'the human condition', just as conflict is not an inherent feature of 'human nature '. Domination and conflict are inherent in class societies, and are based on specific, concrete features of their mode of production. They are rooted in the process of extraction and appropriation of what is produced by human labour. Class domination is not simply a 'fact': it is a process, a continuing endeavour on the part of the dominant class or classes to maintain, strengthen and extend, or defend, their domination.

The focus, always, is on class antagonism and class conflict. This does not mean that Marxism does not recognize the existence of other kinds of conflict within societies and between them ethnic, religious, national, etc. But it does consider these rivalries, conflicts and wars as directly or indirectly derived  from, or related to, class conflicts; whether it is right to do so is not here the point. The fact is that in Marxism this is the essential, primary focus.

Marx himself greatly underestimated and misleadingly belittled his own contribution and that of Engels to this focusing on class antagonism. In a famous letter to his friend, Weydemeyer, dated 5 March, 1852, he wrote that 'no credit is due to me for discovering the existence of classes in modern society, nor yet the struggle between them. Long before me bourgeois historians had described the historical development of this struggle of the classes and bourgeois economists the economic anatomy of the classes'.\endnote{SW 1968, p. 679.} True though this may be (and there is something but no more than something in it), it was nevertheless Marx and Engels who so to speak gave to classes their letters of credit as the dramatis personae of history and to class struggle as its motor; and it is certainly they who, more than anyone else before them, represented politics as the specific articulation of class struggles.

The protagonists of class struggles have naturally varied, through the ages, from 'freeman and slave, patrician and plebeian, lord and serf, guild-master and journeyman', to bourgeoisie and proletariat in the epoch of capitalism. But throughout, 'oppressor and oppressed' have 'stood in constant opposition to one another, carried on an uninterrupted, now hidden, now open fight, a fight that each time ended, either in a revolutionary reconstitution of society at large, or in the common ruin of the contending classes'.\footnote{K. Marx and F. Engels, Manifesto of the Communist Party, in Revs., p. 68. This early allowance by Marx and Engels that one stage of history, rather than inexorably leading to another, may lead to the 'common ruin of all classes', is worth noting.}

The basis of this conflict is somewhat simpler than the endlessly varied forms it assumes in different realms. As already noted, the conflict essentially stems from the determination of the dominant classes to extract as much work as possible from the subject classes; and, conversely, from the attempts of these classes to change the terms and conditions of their subjection, or to end it altogether. In relation to capitalism, the matter is expressed by Marx in terms of the imperative necessity for the owners and controllers of capital to extract the largest possible amount of surplus value from the labour force; and in terms of the latter s attempts either to reduce that amount, or to bring the system to an end. The first alternative involves the attempt to introduce reforms in the operation of capitalism; the second obviously involves its transcendence.

Class domination is economic, political, and cultural - in other words, it has many different and related facets; and the struggle against it is similarly varied and complex. Politics may be the specific expression of that struggle, but, as I noted in the previous chapter, is in fact involved in all its manifestations.

Class domination can never be purely 'economic', or purely cultural': it must always have a strong and pervasive 'political' content, not least because the law is the crystallized form which po itics assumes in providing the necessary sanction and legitimation of all forms of domination. In this sense, 'politics' sanctions what is 'permitted', and therefore 'permits' the rela- tions between members of different and conflicting classes, inside and outside their 'relations of production'.

In the Communist Manifesto, Marx and Engels said that 'the epoch of the bourgeoisie' had 'simplified class antagonisms': Society as a whole is more and more splitting up into two great hostile camps, into two great classes directly facing each other: bourgeoisie and proletariat.'\endnote{Revs., p. 68.} Though the formulation is ambiguous and possibly misleading, it ought not to be taken to mean that capitalist society, which is the 'society'that Marx and Engels are talking about, has been or is being reduced to two classes. It is clear from all of their work that they were perfectly well aware of the continued existence of other classes than the bourgeoisie and the proletariat, and that they did not expect ese classes simply to vanish. Nor is the formulation to be taken to mean that the only antagonism in class society is that between these two classes. They did recognize the existence of other torms of class conflict, and I have already noted that they also recognized the existence of conflicts other than class conflict. The really important point is the insistence by Marx and Engels that the primary conflict in capitalist society is that between capitalists and wage-earners, what Marx, in the formulation w ic quoted in the first chapter, called 'the direct relationship of the owners of the conditions of production to the direct producers'.\endnote{K. Marx, Capital, III, p. 772.} It is this relationship, Marx claimed, which revealed the innermost secret, the hidden basis of the entire social structure';\endnote{Ibid., p. 772.} and which also produced by far the most Class and Class Conflict 21 important element of conflict in capitalist society. Whether this is always the case or not is open to inquiry, and it is probably true that too strict an interpretation of the notion of the primacy of the conflict between capitalists and wage-earners has led to the underestimation by Marxists of the importance which other classes and their conflicts have had and still have in capitalist societies, and to an underestimation also of their role in general --notably that of 'intermediate classes', of which more later.

In political terms, however, the more important question is how far Marx was right in speaking of society as 'splitting up into two great hostile camps'; indeed, whether he was right about this at all. For it would be possible to think of capitalist society 'simplifying' class arrangements without necessarily turning the different classes into mutually hostile ones; or it would at least be possible to think of this hostility, the 'class antagonism' of which Marx speaks without giving it the sharp, warring connotation which he and Engels clearly do give to if in other words to accept that antagonisms will occur between classes and social aggregates of different kinds without producing 'hostile camps', an imagery that significantly conjures up an actual or incipient state of war. This of course raises the whole question of the validity or otherwise of Marx's and Engels's belief in the inevitability of revolution, in other words of a decisive settlement of accounts between bourgeoisie and proletariat, out of which a new social order and mode of production, namely socialism, would emerge.\footnote{In a letter to the Communist Correspondence Committee in Brussels, written from Paris and dated 23 October 1846, Engels reported that, in argument with upholders ol other tendencies, he 'defined the objects of the Communists in this way: (1) to achieve the interests of the proletariat in opposition to those of the bourgeoisie; (2)to do this through the abolition of private property and its replacement by community of goods; (3) to recognize no means of carrying out these objects other than a democratic revolution by force' (SC, p. 37).}

As a preliminary but essential part of the consideration of the many issues which this raises, there is a question which needs to be asked: what is the meaning which is to be attached to the names given to the main or subsidiary antagonists in the struggle, whatever its nature? What, above all, do Marx and Engels, and subsequent Marxist writers, mean when they speak of 'the working class'? The question is obviously crucial. But despite the constant use of the terms in question, or perhaps because of it, their actual meaning is by no means as clear as their use would suggest or than is assumed. Nor indeed is it at all easy to find a ready and obvious answer to the question in the classic writings. In fact, it is quite difficult to find out precisely what Marx meant by the terms 'working class' or 'proletariat'; and later work has not advanced matters very far. The first thing is therefore to try and clarify this and also to identify more closely the other antagonists in the class conflict: this is the basis on which a Marxist politics obviously has to build.

\section*{2}
How genuine and basic the problem is of identifying the exact meaning which is to be attached to the Marxist concept of working class may first of all be gauged from the fact that, for Marx, the very notion of it as a class is in some degree contingent: it is only by fulfilling certain conditions that the working class may properly be said to have become a class. In The Poverty of Philosophy, a passage which refers to the early development of English industrial capitalism makes the point as follows:
\begin{displayquote}
    Economic conditions had first transformed the mass of the people of the country into workers. The domination of capital has created for this mass a common situation, common interests. This mass is thus already a class as against capital, but not yet for itself. In the struggle . . . this mass becomes united, and constitutes itself as a class for itself. The interests it defends become class interests.\endnote{K. Marx, The Poverty of Philosophy, p. 145.}
\end{displayquote}

It is thus in a united struggle, which presupposes a consciousness of its interests, that the proletariat becomes a 'class for itself', as distinct from a mere 'mass' in a common situation and with common interests. In the Communist Manifesto, Marx and Engels went further and spoke of the 'organisation of the proletarians into a class, and consequently into a political party';\endnote{Revs., p. 76. My italics.} and the Manifesto also insists that 'the proletariat during its contest with the bourgeoisie is compelled, by the force of circumstances, to organize itself as a class.'\endnote{Ibid., p. 87.} Rightly or wrongly, a political criterion is thus assigned to the notion of class, and this remained a fundamental theme in Marx's thought. In 1871, the Resolution of the First International on Political Action of the Working Class, drafted by Marx and Engels, still insists that against the 'collective power of the propertied classes the working class cannot act, as a class, except by constituting itself into a political party .\endnote{FI, p. 270.} The ambiguity presented by the formulation whereby the working class constitutes itself as a party will be discussed in Chapter V. But it is worth emphasizing that, for Marx, the working class is not truly a class unless it acquires the capacity to organize itself politically. In so far as this involves will and consciousness, as it obviously does, it can be said that there is in Marx a 'subjective' dimension to the notion of the working class as a class, as well as an 'objective' determination of it. The point might be summarized by saying that, without consciousness, the working class is a mere mass: it becomes a class when it acquires consciousness. What 'consciousness' entails will be discussed presently.

To turn first to this 'objective' determination of the working class: the crucial notion for Marx is that of 'productive worker', which is given its most extensive treatment in Capital and Theories of Surplus Value. The 'productive worker' is he who produces surplus value: 'that labourer alone is productive, who produces surplus-value for the capitalist, and thus works for the self-expansion of capital.'\endnote{K. Marx, Capital (Moscow, 1959), I, p. 509.} Similarly, 'only that wage-labour is productive which produces capital';\endnote{K. Marx, Theories of Surplus Value (Moscow, 1969), Part I, p. 152.} and 'only the wage-labour which creates more value than it costs is productive.'\endnote{Ibid., p. 154.}

It will immediately be seen that this extends the notion of 'worker' far beyond that of the industrial and factory wage- earner. It does so in two ways. Firstly, it covers a large number of people who are not engaged in the industrial process at all, for instance writers, or at least some writers, since, as Marx put it, 'a writer is a productive labourer not in so far as he produces ideas, but in so far as he enriches the publisher who publishes his works, or if he is a wage-labourer for a capitalist.'\endnote{Ibid., p. 158.} In other words, the definition of what constitutes a productive worker has in this conception nothing to do with what he produces: what matters is whether the worker produces surplus value.\endnote{Ibid., p. 157.}

The second extension of the notion of 'productive worker' concerns the actual process of production. Marx writes that

\begin{displayquote}
    the characteristic feature of the capitalist mode of production . . . separates the various kinds of labour from each other, therefore also mental and manual labour - or kinds of labour in which one or the other predominates - and distributes them among different people. This however does not prevent the material product from being the common product of these persons, or their common product embodied in material wealth; any more than on the other hand it prevents or in any way alters the relation of each one of these persons to capital being that of wage-labourer and in this pre-eminent sense being that of a productive labourer.\endnote{Ibid., p. 411.}
\end{displayquote}

In Capital, Marx also succinctly notes that 'in order to labour productively, it is no longer necessary for you to do manual work yourself; enough if you are an organ of the collective labourer, and perform one of its subordinate functions.'\endnote{K. Marx, Capital, I, p. 508. My italics.}

On this view, it is clear that the 'working class' extends far beyond industrial and manual workers. But this extension also creates certain major difficulties for Marxist sociology and politics. For not only does the designation now cover white-collar workers and 'service' workers of every sort, which is not a major difficulty: it also encompasses many other people as well, for instance managerial staff, executive personnel of high rank and even the topmost layers of capitalist production. 'We are all working class now' may be useful conservative propaganda  - but it would be odd to have it legitimated by Marxist concepts; and, more serious, the notion of 'working class' would cease, on this basis, to make possible the differentiations which the class structure of capitalist societies obviously requires.

What is needed here is a set of criteria which do make possible these differentiations - in this instance between the various aggregates of people who constitute different elements of the 'collective labourer', and which allow the necessary distinction to be made between the 'working class' elements of the 'collective labourer' and the rest: between, say, the corporation executive and the factory worker. The criteria in question are partly  - but hardly exhaustively - provided by Marx himself when he refers in the above quotation from Capital to those people who perform the 'subordinate functions' of the 'collective labourer'. The notion of subordination is here crucial, though other criteria of differentiation may be linked to it, for instance income and status, and are usually related to it.

The 'working class' is therefore that part of the 'collective labourer' which produces surplus value, from a position of subordination, at the lower ends of the income scale, and also at the lower ends of what might be called the 'scale of regard'.

This designation does not by any means solve all problems. But neither does any other. One such problem, which is embedded in the notion of class itself, is that of heterogeneity. Like all other classes, the 'working class' is divided into many different strata and by a whole set of differences, which vary according to  time and place, but some of which at least are always present. In the present designation, the main difference would be between industrial wage-earners on the one hand (themselves greatly differentiated) and 'white collar' and 'service' workers on the other and the latter terms obviously cover a wide variation of occupations and grades.

Another such problem, which is a constantly recurring one in the discussion of class, is where to 'cut off' - in this instance to decide at what point (if any) it is appropriate or necessary to draw the line between the 'workers 'just mentioned and the large and growing number of 'workers' who perform a variety of technical, intellectual, supervisory, and managerial tasks. As already noted, these people are indeed part of the 'collective labourer': but whether they are part of the 'working class' is an open question. The point is far from a mere matter of pedantic denomination. On the contrary, it has important political implications, in terms of political strategy and alliances.

As far as classical Marxism is concerned, the 'working class' is basically constituted by industrial wage-earners, factory workers, the 'modern proletariat'. For Marx, Engels, Lenin, and their followers, here is the 'working class', or at least its 'core '. For the purpose of discussing Marxist politics, this will do well enough, provided full account is taken of the many problems which the term presents, precisely in the discussion of politics and such questions as the relation of the working class to its political agencies.

The middle strata of the 'collective labourer' must be distinguished from the so-called 'intermediate' strata of capitalist society of which Marx occasionally spoke\footnote{In Capital, for instance, Marx notes the existence in England of 'middle and intermediate strata which 'obliterate lines of demarcation everywhere' (op. cit. III, p. 862).} and which comprise a wide range of people often also described in Marxist usage as the petty bourgeoisie of capitalist society - medium and small businessmen, shopkeepers, self-employed craftsmen and artisans, small and medium farmers; in other words, that vast and diverse array of people who have not been 'proletarianized', in the sense that they have not become wage and salary earners, and are not therefore part of the 'collective labourer', even though they do of course fulfil definite economic tasks.

In its turn, this petty bourgeoisie must be distinguished from the large and growing army of state employees, engaged in administration and in police and military functions.\footnote{This is not the case for state employees who are engaged in the economic activities covered by the 'public sector', nationalized industries, public services etc., and who obviously are part of the working class.} On the criteria of classification referred to earlier, these state employees are neither part of the working class nor of the petty bourgeoisie: they are, so to speak, a class apart, whose separateness from other classes is bridged by the factor of ideology, which will be considered presently.

To complete this brief enumeration of the main protagonists of class struggle in capitalist society, there remains the capitalist class. This is the class which, for Marx, was so designated by virtue of the fact that it owned and controlled the means of production and of economic activity in general - the great manufacturing, financial and commercial 'interests' of capitalist enterprise. The 'capitalist class', however, extends well beyond these 'interests' and includes many people who fulfil specific professional and other functions on behalf of these 'interests', and who are in various ways - by virtue of income, status, occupation, kinship, etc. - associated with them. It is this variegated totality which is also called the 'ruling class' in Marxist parlance, a concept which needs further discussion.

The point has already been made, but needs to be stressed, that the capitalist class or bourgeoisie (the two terms are used here interchangeably, unless the text requires specific use) is in functional, sociological and in most other terms an heterogeneous class, with many different elements or 'fractions'; and while the development of capitalism has fostered an ever-greater interrelationship between different forms of capital, it has by no means obliterated their differences. There are many issues over which the capitalist class as a whole is more or less united, and this unity may assume a more or less solid political expression, and does assume such expression in times of acute class conflict, 'when the chips are down'. But the economic divisions of the class endure, and so do other divisions of various kinds, according to the particular country in question. The importance of these divisions, from a political point of view, is considerable.

A second question, which has already been encountered in relation to the working class, arises here too, namely the 'cut off' point at which the capitalist class ends and the petty bourgeoisie begins. Marx noted that 'the stratification of classes does not appear in its pure form'\endnote{K. Marx, Capital, III, p. 862.} and this is certainly true here. Clearly a concept of the 'capitalist class' which covers a small entrepreneur, employing half a dozen workmen and the owner of a corporation employing thousands leaves something to be desired. There is no conclusive answer to the problem and some degree of arbitrariness in deciding who, in this context, belongs to the capitalist class is inevitable.

Much more important is the by now well-worn question of ownership and control, or ownership versus control, and the degree to which capitalism and the notion of a capitalist class have been affected by the coming into being of an ever-growing stratum of managers, controlling the most important units of business life, yet doing so without owning more than a minute fraction of the assets they control, and sometimes not even that.

There is no point in rehearsing here the arguments which have been advanced on both sides of the question.\footnote{For a useful recent survey of the argument, see M. Zeitlin, 'Corporate Ownership and ControkThe Large Corporation and the Capitalist Class' in American Journal of Sociology, 1974, vol. 79, no. 5.} My own view of the matter is that managerialism, which had already been noted in its early manifestations by Marx,\footnote{On the basis of the formation of joint stock companies, Marx speaks in Capital of the 'transformation of the actually functioning capitalist into a mere manager, administrator of other people's capital' and of 'private production without the control of private property' (op. cit. Ill, 427, 429).} is indeed a major and growing feature of advanced capitalism; and that the separation of ownership and control which it betokens - when it does betoken it - does not affect in any substantial way the rationale and dynamic of capitalist enterprise. Those who manage it are primarily concerned, whatever they may or may not own, with the maximization of long-term profit and the accumulation of capital for their particular enterprise: ownerless managers are from this point of view practically indistinguishable from owning ones. What matters in both cases are the constraints imposed upon those involved by the imperative and objectively determined requirements of capitalist activity. This being the case, it is perfectly legitimate to speak of a 'capitalist class', occupying the upper rungs of the economic ladder, whatever its members may own, and controlling the operations of capitalist enterprise. It is the more legitimate to do so in that the ideological and political differences between non-owning controllers and the rest have never been more than negligible, if that.

The class struggles in which these classes are engaged occur within the territorial boundaries of the nation-state, and the role of the state in class struggle is of course one of the main objects of attention of Marxist politics. On the other hand, it is as well to stress at the outset that these struggles are deeply and often decisively influenced by external forces. This has been the case throughout the history of capitalism but the point is rendered ever more important by the ever-greater inter-relatedness of capitalism, in a process of 'internationalization' which makes national boundaries an economic anachronism of a constantly more pronounced kind. But the nation state endures, and it is within its boundaries that the encounter occurs.

One or two further preliminary points arise. What encounter? To speak of class conflict is to speak of a central reality by way of metaphor. For classes, as entities, do not enter into conflict - only elements of it do, though it is the case that very large parts of contending classes are on rare occasions directly drawn into battle. For the most part, however, the conflict is fought out between groups of people who are part of a given class, and possibly, though not certainly, representative of it.

Another important question is: What kind of class conflict? For the antagonism between classes does assume many different forms of expression, and many different levels of intensity and scope. It often is strictly localized and focused on immediate, specific and 'economic' demands, and forms part of the 'normal' pattern of relations between employer and wage-earners - with strike action as a familiar part of that pattern. Or it may be fought at the 'cultural' level, and indeed is permanently fought at that level, in so far as there occurs a permanent struggle for the communication of alternative and contradictory ideas, values and perspectives. Or it may be fought at the 'political' level, and bring into question existing political arrangements, large or small. And it may of course assume peaceful or violent forms, and move from one form or level to another.

The distinction between various forms and levels of class conflict is certainly not artificial. But it is nevertheless a misconception to ascribe such labels as 'economic' or 'ideological' to this or that form of conflict. For any event in class conflict, large  or small, includes and expresses all manifestations of social life, and is in this sense an economic, cultural/ideological, social, and political phenomenon. And it is even more important to stress the concomitant proposition, which is a basic part of Marxist perspectives, namely that all manifestations of social life are permanently present in the permanent class conflict of capitalist society.

The class profile which has been outlined here refers to advanced capitalist countries. The very large question which arises is how far that profile and the propositions which relate to it have application to other types of society - the countries of the 'Third World' (using that term with the reservations mentioned earlier) and, even more problematically, the countries of the Communist world. As was noted in Chapter I, the gaps and shortcomings which affect the Marxist theorization of capitalist countries are greatly multiplied in relation to these other forms of society; and I can in any case do no more here than make some cursory remarks about the comparisons and contrasts that may be drawn, in the present context, between these societies and advanced capitalist ones.

In regard to 'Third World' countries, it is clear that class relations are for most of them too the central determinant of their mode of being. But it is equally clear that the classes involved in these relations are in some major ways different or of different importance from those in advanced capitalist societies; and also that, in part because of this and in part for different reasons, the class conflicts engendered by their class relations assume other forms than those encountered in capitalist societies.

The development of these countries has been exceedingly distorted by colonialism and external capitalist domination, direct and indirect; and this has been naturally reflected in their economic, social and political structures. But this also means that Marxism, primarily fashioned in. and for a bourgeois/ capitalist context has, to say the least, to be adapted to the very different circumstances subsumed under the notion of 'underdevelopment'.

One of these different circumstances is that in a large number of these countries, there has existed no strong indigenous class of large-scale capitalists, since the major industrial, extractive, financial, and commercial enterprises are likely to be mainly owned and controlled by foreign interests. The indigenous capitalist class has often tended to be economically rooted in medium- and small-scale enterprise, and partially dependent upon the foreign interests implanted in the country. Correspondingly, the working class is relatively small, compared with the population of the countryside, and concentrated on the one hand in a number of large enterprises and dispersed on the other in a multitude of small ones.

In effect, the mass of the working population is of peasant character, and the main 'relations of production' in these countries tend to be between landlord and peasant, in a multitude of different patterns and connections. But this also means that class conflicts in these economies occur on a very different basis and assume a very different form from those encountered in advanced capitalist countries. This does not mean that Marxist 'guidelines' are inoperative in the analysis of these conflicts. But it does very strongly emphasize the danger of a simple transposition of the Marxist mode of analysis of advanced capitalist societies to countries whose capitalism is of a very different nature.

The same point applies to Communist countries; and it may have to be made with even greater force since there has been a fashion in recent years, strongly encouraged by the Chinese Communists, to claim that the categories of analysis used for capitalist countries would do quite well for the Soviet Union and East European regimes. These, the claim goes, are 'essentially capitalist countries, and the differences between them and other capitalist countries are sufficiently indicated by affixing to them the label 'state capitalist'. Their social structure is one where a 'state bourgeoisie', similar if not identical to the bourgeoisie of Western capitalist countries, exploits and oppresses the working class in the same way as it is oppressed and exploited in classical capitalist countries, and indeed more so. This being so, these regimes are equally susceptible to class conflict, and are only able to contain it by ruthless repression.\footnote{One of the most explicit recent examples of this kind of categorization is C. Bettelheim Les Luttes de classes en URSS (Paris, 1974). For a critical review, see R. Miliband, 'Bettelheim and Soviet Experience', in New Left Review no 91, May-June 1975.}

That there is conflict and repression in these societies is not in question. Nor is the fact of marked disparities of resources, status and power. What is however very questionable indeed is the notion that it is possible to equate these societies with capitalist ones. These are collectivist societies in which the absence of a class which actually owns the means of economic activity, and in which the question of control of these means poses considerable problems, is sufficient to suggest that any such equation is arbitrary and misleading, an exercise in propaganda rather than analysis. Whether the people in these regimes who are at the top of the pyramid are called a class, an elite, a bureaucracy or whatever, it cannot be analysed in the terms which are used to analyse the 'ruling class' of either advanced capitalist societies or of under-developed ones. Nor can the political systems of these collectivist societies. I propose to take up the subject later, and only wish to note here the existence of the problem which it poses for Marxism, and to reiterate the point that it has so far been very inadequately discussed within the Marxist 'problematic '. For the present, the question to which I want to proceed is that of 'class-consciousness' in that 'problematic'.

\section*{3}
Whether or not a class may only be said 'properly'to exist if it has a certain kind of consciousness, it is clear that the element of consciousness is of crucial importance in political terms; and 'class-consciousness' is certainly such an element in Marxist politics. Yet the point has to be made once again that here too there are many more unresolved difficulties than ready usage would suggest.

In Marxist language, class-consciousness may be taken to mean the consciousness which the members of a class have of its 'true' interests - the notion of 'true' interests being one which itself requires elucidation.

The matter is in a sense least complicated in relation to the capitalist class and the bourgeoisie in general. Its true interests presumably consist in the maintenance and defence of capitalism; and its class-consciousness is on this score very easy to achieve. As a matter of historical fact, privileged classes have always been perfectly class-conscious, at least in this sense. On the other hand, the clear perception of the interests of a class in no way betokens a clear perception of the ways in which these interests may best be defended. Also, as a matter of historical fact, privileged classes have very often been short-sighted in this respect, and have needed the skills and adroitness of agents acting on their behalf blit with a sufficient degree of independence to mitigate if not to overcome the short-sightedness of their masters.\endnote{See below, Ch. IV.}

Also, there is in Marxist terms a sense in which the bourgeoisie is falsely conscious, not because it is unable to perceive its true interests, but because it proclaims and believes that these partial and class interests have a universal and classless character. In The German Ideology (1846), Marx and Engels wrote that 'each new class which puts itself in the place of one ruling before it, is compelled, merely in order to carry through its aim, to represent its interest as the common interest of all the members of society, that is, expressed in ideal form: it has to give its ideas the form of universality, and represent them as the only rational, universally valid ones'.\endnote{K. Marx and F.Engels, The German Ideology, p. 63.} In writing thus, Marx and Engels were thinking primarily of the bourgeoisie's struggles against feudal rule, and particularly of the French bourgeoisie's protracted struggles for intellectual as well as economic and political ascendancy under the ancien regime. 'Ideology', for Marx and Engels, is precisely the attempt to 'universalize'and give ideal form to what are no more than limited, class-bound ideas and interests: it is in this sense that they use the word 'ideology' pejoratively, as meaning a false representation of reality. At the same time, they did not hold to the vulgar view that this false representation was necessarily deliberate. Deliberate deception does indeed occur, whereby the spokesmen of a dominant class act as the ideologues' of that class, and try to peisuade the subordinate classes of the universal validity of ideas and principles which these spokesmen know to be partial and class-bound but useful in the maintenance of the given social order. But alongside deliberate deception, there is also much, and perhaps more, of self-deception, in so far as the spokesmen of a dominant class, and those for whom they speak, do deeply believe in the universal truth of the ideas and ideals which they uphold, and will therefore fight for them all the more vigorously and, if need be, ferociously.

The obvious question which arises here is why the point should not also apply to the working class - in other words why should the working class be thought of and claim to be a 'universal class, whose interests are indistinguishable from those of  society at large, which Marx and Engels did claim to be the case, and which has remained a central Marxist claim ever since.

The answer which they gave to that question has many facets and ramifications which need to be examined and followed through, but the basis of it is that the working class is not only the vast majority of the population but that it is also the only class in history whose interests and well-being do not depend on the oppression and exploitation of other classes. 'All previous historical movements', they said in the Communist Manifesto, 'were movements of minorities, or in the interest of minorities.

The proletarian movement is the self-conscious, independent movement of the immense majority, in the interest of the immense majority'.\endnote{Revs., p. 78.} Indeed, the proletariat, having swept away the old conditions of production, 'will, along with these conditions, have swept away the conditions for the existence of class antagonisms and of classes generally, and will thereby have abolished its own supremacy as a class'.\endnote{Ibid., p. 87.} All previous revolutions had by necessity been limited in scope, because of the narrow class interests of those who had made them. By contrast, 'the communist revolution is the most radical rupture with traditional property relations'; and, Marx adds, 'no wonder that its development involves the most radical rupture with traditional ideas'.\endnote{Ibid., p. 86.} Related to this, and reinforcing the notion of the working class as a 'universal ' class, there is the Marxist view that the working class alone is capable of acting on behalf of the whole of society and remove from it the greatest and weightiest of all impediments to its boundless development, namely the capitalist mode of production. How far some of these claims are themselves tainted with 'ideology' is an interesting and important question.

What then is class-consciousness in reference to the working class - the class with which Marxists in this context have been mainly concerned? A great deal, in terms of Marxist politics, hinges on the answer to this.

In the Marxist perspective, proletarian class-consciousness may be taken to mean the achievement of an understanding that the emancipation of the proletariat and the liberation of society require the overthrow of capitalism; and this understanding may also be taken to entail the will to overthrow it. It is in this sense that proletarian class-consciousness is also revolutionary consciousness.

It is of course possible to invest the notion of revolutionary consciousness with any meaning one chooses; and to argue that no one can be said to be 'truly' or 'really' class-conscious who does not subscribe to this or that idea, precept, strategy, party, and whatever else. But it is essential to realize that Marx himself did not define it with any such specificity, and that it did remain for him, as for Engels, a general concept, without any more precise meaning than that suggested above. In fact, it is very striking that Marx and Engels consistently and vigorously dismissed the notion that there was a set of ideas which specifically defined revolutionary consciousness. In the Communist Manifesto, they said that the Communists 'do not set up any sectarian principles of their own, by which to shape and mould the proletarian movement ';\endnote{Ibid., p. 79.} and they also insisted that 'the theoretical conclusions of the Communists are in no way based on ideas or principles that have been invented, or discovered, by this or that would-be reformer'.\footnote{Ibid., p. 80. In his Class Struggles in France (1850), Marx contemptuously referred to 'petty-bourgeois socialists' for whom 'the coming historical process' appeared 'as an application of systems, which the thinkers of society, either in company with others, or as single inventors, devise or have devised. In this way they become the eclectics or adepts of exi sting socialist systems, of doctrinaire socialism, which was the theoretical expression of the proletariat only as long as it had not yet developed further and become a free, autonomous, historical movement' (SE, p.122). The italics, it will be recalled, are in the text.} Many years later, Marx wrote in The Civil War in France that the working class 'have no ready-made utopias to introduce par decret du peuple' and even that 'they have no ideals to realize, but to set free the elements of the new society with which old collapsing bourgeois society itself is pregnant'.\endnote{FI, p. 213.} Soon after, Marx and Engels, in a Circular concerning 'The Alleged Splits in the International' (1872), noted that, while the rules of the International gave to its constituent societies a common object and programme, that programme was 'limited to outlining the major features of the proletarian movement, and leaving the details of theory to be worked out as inspired by the demands of the practical struggle, and as growing out of the exchange of ideas among the sections, with an equal hearing given to all socialist views in their journals and congresses'.\footnote{Ibid., p. 299. Note also Marx's warning in a private letter to his daughter and son-in-lawinl870:'Sectarian "etiquettes" must be avoided in the International. The general aims and tendencies of the working class arise from the general conditions in which it finds itself. Therefore, these aims and tendencies are found in the whole class, although the movement is reflected in their heads in the most varied forms, more or less imaginary, more or less related to the conditions. Those who best understand the hidden meaning of the class struggle which is unfolding before our eyes - the Communists - are the last to commit the error of approving or furthering Sectarianism ' (K. Marx to Paul and Laura Lafargue, in K. Marx and F. Engels, Werke (Berlin, 1965), vol. XXXII, p. 671).}

The general point is one of extreme importance. In relation to revolutionary consciousness, it means that, for Marx and Engels at least, the concept did not entail the kind of devout and categorical adherence to given formulas which became the hallmark of later Marxism. Indeed, they never even claimed that the achievement of class-consciousness required adherence to something specifically called Marxism - Marx's cast of mind suggests that he would have found the idea rather laughable\footnote{The case has occasionally been made that it was Engels who, after Marx's death, first 'invented' Marxism as a political creed and as the political doctrine of the working-class movement. For an extreme but scholarly statement of this thesis, see M. Rubel.Marxcritiquedu marxisme (Paris, 1974), particularly Ch. I.}.

Unfortunately, later developments turned the issue into anything but a laughing matter. For the attempt was made from many quarters, and has not ceased to be made, to stipulate exactly what revolutionary consciousness means in terms of convictions on a vast range of questions; and also what convictions on an equally vast range of questions such revolutionary consciousness must exclude. One consequence of this is to turn class-consciousness into a catechismal orthodoxy, departure and dissent from which become grave - and punishable - offences. Another consequence is to enhance tremendously the role of the keepers of the orthodoxy, namely the party leaders and their appointees: if revolutionary consciousness can be so precisely defined, there must be an authority to define it, and to decide when and in what ways it must be modified.

Even the notion of class-consciousness as something to be 'achieved ' is not free from question-begging connotations. For it suggests a state to be reached, a thing to be appropriated; or at least this is what it may well imply. But in this meaning, it robs the concept of its dynamic nature, deprives it of its character as a process, which is constantly changing, and which is susceptible not only to advance but to regress, and certainly in no way unilinear. In other words, revolutionary consciousness is not some kind of Marxist state of grace which, once achieved, is total and irreversible. It is a certain understanding of the nature of the
social order and of what needs to be done about it. As such, it comprises many uncertainties, tensions, contradictions, open questions, and possibilities of error and regression. This at least is how it has always been in reality.

These qualifications do not, however, deprive the notion of class\hyp{}consciousness, in its Marxist sense, of considerable and distinctive meaning. For it does after all denote a commitment to the revolutionary transformation of society, an 'interiorization' of the need to achieve that 'most radical rupture' with traditional property relations and traditional ideas of which Marx and Engels spoke in the Communist Manifesto and to which they held throughout. This meaning of class\hyp{}consciousness does not carry in its train specific answers to many questions of great importance to the working\hyp{}class movement. But it does carry certain definite perspectives and establishes certain firm delimitations. No more - but no less either.

The class with whose class-consciousness Marxists, beginning with Marx and Engels, have always been preoccupied is of course the working class; and the latter's relation, so to speak, to class-consciousness needs to be considered. But before doing so, it is worth noting that it is not only the proletariat which, in Marxist terms, can achieve class-consciousness. Again in the Communist Manifesto, Marx and Engels said that

\begin{displayquote}
    in times when the class struggle nears the decisive hour, the process of dissolution going on within the ruling class, in fact within the whole range of old society, assumes such a violent, glaring character, that a small section of the ruling class cuts itself adrift, and joins the revolutionary class, the class that holds the future in its hands ... a portion of the bourgeoisie goes over to the proletariat, and in particular, a portion of the bourgeois ideologists, who have raised themselves to the level of comprehending theoretically the historical movement as a whole.\endnote{Revs., p. 77.}
\end{displayquote}

Half a century later, Lenin echoed this in relation to 'educated representatives of the propertied classes, the intellectuals', and noted that 'the founders of modern scientific socialism, Marx and Engels, themselves belonged to the bourgeois intelligentsia.\endnote{What is to be Done? in CWL, vol. 5 (1961), p. 375 27 EW, p. 256.}

Such instances of class betrayal on the part of members of the bourgeoisie, whether intellectuals or not, have been quite common in the history of the working -class movements; and many of  the leaders of these movements have been of bourgeois origin. But for the present discussion, the more important question is that posed by the class orientations of other classes, notably the middle strata of the 'collective labourer' referred to earlier; and of the petty bourgeoisie.

About the latter, Marx and Engels were generally scathing, and their view of the petty bourgeoisie has been endorsed by later Marxists - and for that matter confirmed by historical developments. 'The lower middle class, the small manufacturer, the shopkeeper, the artisan, the peasant, all these fight against the bourgeoisie', said Marx and Engels, but only 'to save from extinction their existence as fractions of the middle class'; and this meant that 'they are therefore not revolutionary, but conservative. Nay more, they are reactionary for they try to roll back the wheel of history.\footnote{Revs., p. 77. Marx and Engels were equally or even more scathing about the 'lumpenproletariat', the 'social scum, that passively rotting mass thrown off by the lowest layers of society', most of whom could be expected to become 'the bribed tool of reactionary intrigue' (ibid., p. 77).}

Summary though the characterization may be, it does quite accurately pinpoint the general position adopted by members of the petty bourgeoisies of advanced capitalism - certainly most of them have been fierce opponents of organized labour, and unwilling allies, but allies none the less, of the large-scale capitalist interests which threaten their economic existence. This 'intermediate' class may be wooed and even pacified by a workers' movement, but its members can hardly be expected fundamentally to change a class-consciousness which is deeply rooted in their economic and social circumstances, and which sets them at odds with proletarian class-consciousness.

The issue presents itself very differently in regard to the vast and ever growing number of people who man the technical, scientific, supervisory, and cultural posts of advanced capitalist societies. This 'new working class', as it has been called, is sharply pulled in contrary directions. On a variety of economic, social, and cultural criteria, it is markedly differentiated from the 'traditional' working class; and some of its members may have plausible hopes of access to the upper layers of the social pyramid. On the other hand, it has undergone a steady process of 'proletarianization', in so far as it is a salaried and subordinate part of the 'collective worker': and it has in recent decades learnt the virtues of collective organization and collective action in defence of its sectional interests. Such barriers as do exist to its development of 'class-consciousness' are not insuperable, or for that matter particularly high; and this may have very large repercussions indeed on the political plane, since an organic linkage between this part of the 'collective worker' and the rest of the working class is bound to make a considerable difference to the nature and impact of the political organizations of the working class. This is one of the most important open areas of the political sociology of Marxism, and of its politics.

But the really big question still remains: why should the working class have or acquire revolutionary class-consciousness the understanding that it must do away with capitalism to emancipate itself and society ? Why should it not reject the call to revolution, and seek reform of various kinds within the loose confines of capitalism, in accordance with what Lenin and others after him described as mere 'trade union consciousness ? These questions are the more relevant since the working class has, or so it is claimed, generally chosen this second path and resolutely refused to act out the revolutionary role assigned to it by Marxism.

An early answer to this kind of question was given by Marx in his introduction to the Critique of Hegel's 'Philosophy of Right'. Discussing the inability of the German bourgeoisie to make a thoroughgoing revolution, Marx asked: 'Where is the positive possibility of German emancipation ?'; and he goes on to say that the answer lies 'in the formation of a class with radical chains, a class of civil society which is not a class of civil society, a class which is the dissolution of all classes, a sphere which has a universal character because of its universal suffering and which lays claim to no particular right because the wrong it suffers is not a particular wrong but wrong in general\dots'\endnote{EW, p. 256.}

That class, he said, was the proletariat. As if to anticipate later developments and objections, Marx and Engels also wrote in The Holy Family (1844) that 'the question is not what this or that proletarian, or even the whole of the proletariat at the moment considers as its aim. The question is what the proletariat is, and what, consequent on that being, it will be compelled to do. Its aim and historical action is irrevocable and obviously demonstrated in its own life situation as well as in the whole organisation of bourgeois society today.'\endnote{K. Marx and F. Engels, The Holy Family (Moscow, 1956), p. 53 29 FI, p. 79.}

Marx's early formulations of the 'role ' and the 'mission ' of the proletariat as an agent of emancipation do undoubtedly have a fairly heavy Hegelian imprint, with the proletariat almost occupying in the unfolding of history the role which Hegel assigned to the Idea. But even in these early formulations, there is in Marx and Engels a concept of the proletariat as destined to become a revolutionary class because revolution is its only means of deliverance from the oppression, exploitation and alienation which existing society imposes upon it. These features of existing society are inherent to it, an intrinsic part of this social order, and can therefore only be got rid of by the disappearance of the social order itself. From this point of view, the proletariat's role is not determined by any extra-historical agency: it is determined by the nature of capitalism and by the concrete conditions which it imposes upon the working class and upon society at large.

Against this, it has often been argued that conditions have greatly changed over time, and that Marxists, beginning with Marx, have always had too strong a tendency to under\hyp{}estimate the capacity of capitalism to assimilate far\hyp{}reaching reforms in every area of life.

However true this may be, it misses the real argument, which is that capitalism, however many and varied the reforms it can assimilate, is unable to do without exploitation, oppression, and dehumanization; and that it cannot create the truly human environment for which it has itself produced the material conditions.

This, on the other hand, leaves unanswered the objection that the working class has not, over a period of a hundred years and more, developed the class-consciousness which Marxists have expected of it and turned itself into a revolutionary class. This needs to be looked at more closely, since it relates to the meaning of class-consciousness as it concerns the working class.

One fairly common element of confusion in this discussion is the equation of revolutionary consciousness with the will to insurrection - the absence of such a will in the working class being automatically deemed to demonstrate a lack of class- consciousness. But this is precisely the kind of arbitrary designation of what does not specifically constitute such consciousness, which Marx never sought to lay down. It may well be that class-consciousness and revolutionary consciousness must eventually come to encompass a will to insurrection; and Marx himself, without being dogmatic about it, and while allowing for some possible exceptions, did believe that the abolition of capitalism would require its violent overthrow. But the will to insurrection which this would entail must be seen as the ultimate extension of revolutionary consciousness, as its final strategic manifestation, produced by specific and for the most part unforeseeable circumstances. That the working class has only seldom, and in some countries never, manifested much of a will to insurrection is not in Marxist terms a decisive demonstration of a lack of class-consciousness.

Nor for that matter is the pursuit of specific and partial reforms within the ambit of capitalism and through the constitutional and political mechanisms of bourgeois regimes. Marx himself vigorously supported the pursuit of such reforms and it may be recalled that, in the Inaugural Address of the First International, which he wrote in 1864, he hailed one such reform, the Ten Hours Bill, as 'not only a great practical success' but as 'the victory of a principle; it was the first time that in broad daylight the political economy of the middle class succumbed to the political economy of the working class.\endnote{FI, p. 79.} Nor did Engels have the slightest difficulty in giving his support to the parliamentary, electoral, and quite 'reformist' endeavours of the German Social Democratic Party, notwithstanding certain reservations - and even these should not be exaggerated.\endnote{See below, pp. 79-80.} As for Lenin, his whole work is permeated by firm approval for the struggle for partial reforms of every sort, including the most modest 'economic' reforms; and so is his work peppered with contemptuous denunciations of the all-or-nothing approach to the revolutionary struggle, culminating in 'Left-Wing' Communism - An Infantile Disorder of 1920.

The question is not one of support for reforms: to designate such support as an example of 'false consciousness' is as arbitrary as the equation of 'false consciousness ' with the absence of a will to insurrection. The real issue is the perspective from which reforms are viewed, what they are expected to achieve, and what else than reforms is being pursued. What Lenin called 'trade union consciousness', and which he counterposed to revolutionary consciousness was a perspective which did not go beyond the achievement of partial reforms and was content to seek the amelioration and not the abolition of capitalism.

Marxists have always believed that the working class would eventually want to go beyond partial reforms inside the system - that it would, in other words, come to acquire the 'class- consciousness' needed to want a thoroughgoing, revolutionary transformation of capitalist society into an entirely differently based and differently motivated system.

In some meanings of 'class-consciousness ' and of 'revolutionary consciousness', this hope might well amount to what C. Wright Mills scathingly called a 'labour metaphysic', grounded in nothing more than faith. But in the meaning which it had for Marx and Engels, this is not what it amounts to. It simply derives from the conviction that the working class, confronted with the shortcomings, depredations, and contradictions of capitalism, would want to get rid of it, in favour of a system which, as Marx put it, transformed 'the means of production, land and capital, now chiefly the means of enslaving and exploiting labour, into mere instruments of free and associated labour'.\endnote{K. Marx, The Civil War in France, in FI, p. 213.}

These propositions leave open many questions, but there is no 'labour metaphysic' about it; so little, indeed, in the case of Lenin, that he bluntly said in 1902 that 'the working class, exclusively by its own effort, is able to develop only trade-union consciousness'.\endnote{V. I. Lenin, What is to be Done?, op. cit., p. 375.} This too has very large implications for Marxist politics. But it may be said at this stage that a strict reading of the record, both of capitalism and of the working class, suggests that the expectation of a growing development in the working class (and in other large segments of the 'collective labourer') of a will for radical change is not in the least unreasonable or 'metaphysical'.

Neither Marx nor any of the classical Marxist writers had any illusion about the massive obstacles which the working class would have to overcome on the way to acquiring this class- consciousness, and about the difficulties that there would be in breaking through the fog of what Gramsci called the 'common- sense' of the epoch. They knew well, as Marx put it in The Eighteenth Brumaire of Louis Bonaparte of 1852 that 'the tradition of the dead generations weighs like a nightmare on the minds of the living ';\endnote{SE, p. 146.} and that enormous efforts would be undertaken by those in whose interest it was to do so to make the weight heavier still. I now propose to discuss the nature of the obstacles which obstruct and retard the acquisition of class- consciousness by the working class.

\theendnotes